
%\documentclass[oneside,final,14pt]{extreport}
\documentclass[12pt, a4paper, openany]{book}
\usepackage[left=2cm,right=2cm,top=1.5cm,bottom=1.5cm,bindingoffset=0cm]{geometry}
%\usepackage[koi8-r]{inputenc}
\usepackage[russianb]{babel}
\usepackage{vmargin}
\setpapersize{A4}
\usepackage[T2A]{fontenc}
\usepackage[utf8x]{inputenc} % more recent versions (at least>=2004-17-10)
%\usepackage[russian]{babel}
%\setmarginsrb{2cm}{1.5cm}{1cm}{1.5cm}{0pt}{0mm}{0pt}{13mm}
\usepackage{indentfirst}
\usepackage{nicefrac} % For comparison
%\usepackage{xfrac} % Works better with other fonts
%\usepackage[unicode, pdftex]{hyperref}
\usepackage{lettrine}
\usepackage[usenames]{color}
%\special{papersize=a5}
\usepackage{colortbl}
\usepackage[pagestyles]{titlesec}
\usepackage{xfrac} % Works better with other fonts
\usepackage[colorlinks=true,linkcolor=black,urlcolor=black,bookmarksopen=true]{hyperref}


% Настройка вертикальных и горизонтальных отступов
\titlespacing{\chapter}{0pt}{5pt}{5pt}
\titlespacing{\section}{\parindent}{4mm}{4mm}
\titlespacing{\subsection}{\parindent}{3mm}{3mm}

\newcommand{\anonsection}[1]{ \section*{#1} \addcontentsline{toc}{section}{\numberline {}#1}} 

\makeatletter %%%%% <---- Starting chapter without a pagebreak
\renewcommand\chapter{\par%
	\thispagestyle{plain}% \global\@topnum\z@
	\@afterindentfalse \secdef\@chapter\@schapter}
\makeatother %%%%% <---- Starting chapter without a pagebreak
\titleformat{\chapter}[display]
{\normalfont\bfseries}{}{0pt}{\Large}

\newpagestyle{mystyle}{
	\sethead[\thepage][][]{}{}{\thepage}	
}

\pagestyle{mystyle}

\sloppy
\begin{document}

	
	\begin{titlepage}
		
		\begin{center}

			\vspace*{\fill}
			
			
			{\large\bf ЛЮБОВЬ АГЕЕВА\\}
			\ \\
			\ \\
			{\Huge\bf КАЗАНСКИЙ ФЕНОМЕН: МИФ И РЕАЛЬНОСТЬ\\}
			\ \\
			\ \\
			\ \\
			\vspace*{\fill}
			
			\vfill
			

			
			КАЗАНЬ
			
			ТАТАРСКОЕ КНИЖНОЕ ИЗДАТЕЛЬСТВО
			
			1991

		\end{center}
		
	\end{titlepage}
	
	\thispagestyle{empty} % выключаем отображение номера для этой страницы
	
	\newpage
	

\setcounter{secnumdepth}{0}

	\section[Несколько слов вместо напутствия]{\center \textbf{НЕСКОЛЬКО СЛОВ ВМЕСТО НАПУТСТВИЯ}}
	
	\newpage
	
	\section[Вчера это было в наших силах]{\center \textbf{ВЧЕРА ЭТО БЫЛО В НАШИХ СИЛАХ}}
	
	\subsection[Предыстория]{\center \textbf{ПРЕДЫСТОРИЯ}}

Листаю старые подшивки своей газеты и пытаюсь понять, почему журналисты «проморгали» это явление.

На школьные темы «Вечерка» писала много. Осо­бенно последовательны мы были в освещении проблем трудового воспитания и профессиональной ориентации. Активно провели обсуждение проекта Основных направ­лений реформы общеобразовательной и профессиональ­ной школы. Писем от читателей пришло много, с кон­кретными предложениями и советами. О драках раз­говора не было.

Город еще спокоен...

Рассказывая 4 сентября 1980 года о суде над под­ростком Леонидом С., я не могла предположить, что за этим конкретным уголовным делом сокрыта боль­шая опасность. Подросток попал на скамью подсуди­мых как активный участник драки. Ватага его сверст­ников пошла в соседний поселок, чтобы выяснить отно­шения с тамошними мальчишками. На учет в инспекции по делам несовершеннолетних Леня был поставлен го­дом раньше —за такую же драку.

Что это было — предчувствие беды или уже сама беда? Ведь в апреле 1980 года Верховный суд ТАССР рассмотрел знаменитое дело банды «Тяп-ляп».

В 1983 году первокурсник отделения журналисти­ки КГУ Артур Гафаров опубликовал в «Комсомольце Татарии» статью о дворовой компании с оригинальным названием «Грязь». Статья называлась «Кого ты хотел удивить?» и была написана по собственным впечатлениям. Артур был членом этой компании, а когда захотел из нее выйти, оказался не раз битым. Потом у него была еще одна публикация на эту тему — «Престиж улицы подлинный и мнимый». Но мимо не придал значения публикациям студента.

В нашей газете первым поведал о группировках корреспондент Н. Алкин — это было в том же 1983 году. Рассказывая о том, как на хоккейной площадке школы №136 избили Володю Н., он с позиции взрослого человека обобщил свои впечатления о подростках, о судебном процессе.

В 1984 году корреспондент газеты «Советская Татария» С. Кулешова писала об опасных игрушках подростков: обыкновенные баллончики из-под бытовых сифонов они переделывали в «бомбочки», что нередко заканчивалось травмой.

В нашей газете после большого перерыва, в январе 1986 года, выходит  корреспонденция Г. Наумова «Шел по городу троллейбус». Если об избиении подростка во дворе школы №135 знали до газетной публикации немногие, то о разгромленном троллейбусе говорил весь город.

И все-таки две разрозненные статьи плюс публи­кации в других газетах спокойствия горожан не нару­шили. На пресс-конференции председателя горисполкома Р. Идиатуллина, организованной нашей редакцией 21 февраля 1986 года, не  было задано ни одного вопроса о подростках. На собрании партийно-хозяйствен­ного актива города, состоявшегося в октябре месяце, о драках не было произнесено ни слова.

Судя по нашим публикациям, должностные лица уже искали выход из положения. Например, в июне «Вечерка» опубликовала проект решения будущей сессии  горсовета, на которой планировалось обсудить ме­роприятия по организации досуга подростков и моло­дежи.

Город еще спокоен, но проявляем беспокойство мы, журналисты. На столе у редактора «Вечерней Ка­зани» А. Гаврилова лежат четыре мои статьи с ана­лизом ситуации (1986: «Словно на разных языках», 6 июня; «Воспитание наказанием», 7 июня; «Этот неслучайный случай», 8 июня; «Спроси с себя», 9 июня), но горком КПСС, точнее секретарь по идеологии Ф. Зиятдинова, не рекомендует их публико­вать. Это тем более странно, что статьи были написа­ны по ее просьбе.

Не спасла высокая оценка профес­сионала: доцент Казанского филиала Высшей юриди­ческой заочной школы Н. Фатхуллин в рецензии для внут­реннего, горкомовского пользования писал: «Статьи от­личаются глубиной анализа, целостным концептуаль­ным подходом к решению сложной и актуальной проб­лемы, пронизаны озабоченностью за судьбы подрастаю­щего поколения. Вместе с тем это трезвый, по-партий­ному принципиальный взгляд на реальную действитель­ность, смелое выявление подлинной природы негативных тенденций, в том числе участившихся правонару­шений несовершеннолетних. Последовательно объектив­ный анализ обеспечивает и реальность выводов, на­дежность рекомендаций».

Эту рецензию Нариман Садреевич впоследствии, когда мы будем работать над проблемой сообща, пода­рит мне на память.

Только после вмешательства первого секретаря горкома партии М. Сухова мои статьи были опубликова­ны — спустя несколько месяцев после написания. После этого рубрика «Город и подросток» с полос газеты уже не сходила.

Чем же объяснить отсутствие публикаций о вражде группировок? Сегодня можно прикрываться запретами «сверху». Однако запретов на подростковую тему в пору становления грозного явления не было. Во вся­ком случае, для журналистов нашей газеты. Для нас, как и для всех, существовали лишь единичные факты, которым мы не придали серьезного значения. Прихо­дится самокритично это признать. Правда, мы на сей счет обладали очень малой информацией. На учительских совещаниях, где я обязательно присутствовала, о драках не говорилось, на «узкие», где они обсуждались, нас не приглашали.

Видимо, сначала дракам не придавали особого значения, а когда ситуация стала обостряться, решили, по давней привычке, сор из избы не выносить, справиться  собственными силами, без журналистов и общественности. Когда ситуация вышла из-под контроля и в Ка­зань зачастили различные комиссии из Москвы, прибегли к прямым запретам. Тогда это еще было возможно, как в случае с моими статьями. Но на другой запрет — не упоминать в газете слово «дискотека» — мы уже не обратили внимания. Даже из горкомовского кабинета нельзя было запретить новую форму молодежного досуга.

Позиция страуса нам дорого стоила. Время для анализа и уж тем более активного вмешательства в ход событий было упущено. Группировки стали настолько сильны, традиции их вражды столь живучи, что не было никаких надежд на успех с первой кавалерийской атаки. Правда, кое-кто на это упо­вал.

Общественность постоянно успокаивали: группировок в городе не более десятка, они малочисленны и опасности не представляют; большинство подростков — нормальные ребята... Так болезнь загонялась внутрь, желаемое выдавалось за действительное. Старая бо­лезнь периода застоя. И хотя на дворе уже была перестройка, болезнь эта оставалась.

Сегодня, с высоты накопленного опыта, видны мно­гие причины, послужившие детонатором явления, которое потом получит название «казанский феномен». Мои публикации того времени говорили о деформациях личности   ребенка, подростка, об отчуждении его от семьи, от школы. Как это ни горько сознавать, мы сами толкали ребят в группировки.

Анализируя тогда конкретные судьбы, я, как пра­вило, не выходила на какие-то серьезные обобщения, касающиеся социальных причин деформаций. В то время считали, что таких причин просто не существует. Были отдельные недоработки отдельных родителей и отдель­ных учителей, некрасивые проступки отдельных ре­бят.

Но размышляя над судьбой своих героев, я все чаще натыкалась на одно явление, о котором однажды рискнула написать в «Литературную газету». Мне дове­лось видеть гранки своей статьи с крупным заголов­ком «Показуха», однако читатели с ней так и не познакомились. Меня попросили написать другой ма­териал, вышедший в свет в скромном виде читатель­ского письма.

Корреспондент «ЛГ» Елена Шиманко объяснила мне, что статья не понравилась высокому должностно­му лицу из ЦК КПСС, имевшему право вето на газетные материалы. Он не без оснований заметил, что после очередного съезда ВЛКСМ, прошедшего, как всегда, с большой помпой, у читателей могут возникнуть нежелатель­ные ассоциации.

	\newpage

\section[Словарь казанских подростков]{\center \textbf{СЛОВАРЬ КАЗАНСКИХ ПОДРОСТКОВ}}
Молодежная среда всегда имеет свой жаргон. Какие-то словечки живут в ней годами. Например, предупредительный клич об опасности «атас», который я помню с детства, встречается и сегодня. Есть понятия, связанные с определенным временем, с определенным уровнем общественного развития. Наконец, возможны какие-то новообразования, характерные для определенного региона страны. 

Казань в этом отношении являет пример показательный. Новое социальное явление в среде молодежи не могло не оставить своего следа в языке. Словарь казанских подростков составляют достаточно много понятий, которые определяют ранее не существующие отношения. Одни из них – «мотаться», «мотальщики»– стали широко известны благодаря публикациям центральных газет и телепередачам, другие обитают в пределах Казани и Татарии.

Чтобы читатели лучше смогли ориентироваться в жизни казанских подростков, предлагаю вашему вниманию словарь наиболее употребительных жаргонизмов, введенных в русский язык группировками.

\textbf{Мотаться} – быть членом группировки, жить по ее законам.

\textbf{Мотальщики, мотаны, гопники} – так называют членов группировок.

\textbf{Контора} – молодежная компания, объединяющая молодежь определенного двора, микрорайона, улицы. На официальном языке конторы называют группировками.

\textbf{Шелуха, салаги} – члены группировки 10–13 лет. Это резерв «конторы».

\textbf{Супера} – члены группировки 14–16 лет.

\textbf{Молодые} – члены группировок 17–19 лет, входящие в руководящее ядро молодежных компаний.

\textbf{Старики, старшие} – юноши и взрослые люди. Возрастные границы – от 20 до 35 лет и даже старше.

\textbf{Авторитеты} – высшее звено руководства группировкой. Их еще называют «стариками», «дедами», «королям».

\textbf{Автор} – руководитель определенного возрастного отряда, определяющий жизнедеятельность группы. Он входит в руководящее ядро группировки, обеспечивает связь своего отряда с другими частями, компаниями.

\textbf{Коробка} – исходное ядро группировки, которое включает двор или несколько рядом стоящих многоэтажных домов. Группировку может образовать одна «коробка», но есть «конторы», включающие несколько «коробок».

\textbf{Сбор, обязаловка} – регулярные встречи членов группировки в точно назначенные дни и часы.

\textbf{Разборка} – обсуждение всей группировкой «боевых действий» – набегов на «чужую территорию», драк...

\textbf{Чушпан} – подросток, не состоящий в группировке, тот, кто не мотается.

\textbf{Копилка} – черепно-мозговая травма.

\textbf{Грязнушка (Грязь), Жилка, Чайники, Пентагон, Двадцатый двор, Воровский, Финны, Стандартный поселок} и другие – названия группировок.

\textbf{Общак} – коллективная касса группировки, которая состоит из добровольных взносов, разовых и регулярных, а также из поступлений, добытых преступным путем.

\textbf{Общая девочка} – при многих группировках есть девичьи компании, полностью разделяющие законы их жизни. Они вместе проводят время, являются сексуальными партнершами всей компании.  	
	
	\newpage
	\tableofcontents
	
	\thispagestyle{empty} % 
	
	\newpage
	
	\setcounter{secnumdepth}{0}
	
	\phantomsection
	
		\section*{Описание}
	
	{\bf Название:} Казанский феномен: миф и реальность
	
{\bf Автор:} Любовь Владимировна Агеева
	
{\bf Издательство:} Казань: Татарское кн. изд-во.
	
		{\bf Аннотация:} Автор сборника пытается проанализировать подростковые проблемы города, осмыслить изнутри так называемый <<казанский феномен>>. Исследую причины возникновения подростковых группировок, журналисты, ученые, родители высказывают разные точки зрения, знакомят с различными подходами в решении этой проблемы.
		
		Адресуется широкому кругу читателей.
		\thispagestyle{empty} % выключаем отображение номера для этой страницы

	
\end{document}


